\section{Basin geometry of grokking initializations}
\label{sec:basin-geometry}

The previous notes describe \emph{what} the trained model implements (the
Fourier multiplication algorithm of \cref{sec:mod-mul-algorithm}) and how that
algorithm can be implemented unusually efficiently when capacity is tight
(\cref{sec:legendre-channel}).  Both are statements about the end state of
training.  This note begins a third thread: \emph{which} end state is reached
from a given initialization, and what dynamics during training determine that
selection.

Empirically we have seen that under identical hyperparameters at
$d_{\text{model}} = 24$, different random seeds find different $\mathcal K$,
and one seed in eight fails to grok cleanly.  We characterize the set of
``good'' initializations and the dynamical mechanism by which they reach
their algorithmic basins.

\subsection{Setup}
\label{subsec:bgeo-setup}

We work with the architecture and training recipe of \cref{subsec:setup}: a
one-layer transformer, no LayerNorm, ReLU MLP of width $d_{\text{mlp}}$,
trained by AdamW on full-batch cross-entropy with weight decay $\lambda > 0$.
The embedding matrix is $\WE \in \R^{d_{\text{model}} \times p}$; only its
restriction $W := \WE[:, 0\!:\!p]$ to value-token columns is dynamically
relevant for the algorithm.

We approximate the optimizer by gradient flow with weight decay:
\begin{equation}
  \frac{\mathrm dW}{\mathrm d t} \;=\; -\nabla_W \mathcal L(W; \theta_{\text{rest}}) \;-\; \lambda W,
  \label{eq:bgeo-gradflow}
\end{equation}
where $\theta_{\text{rest}}$ are the remaining weights.  The cross-entropy loss is
\begin{equation}
  \mathcal L \;=\; -\frac{1}{|\mathcal D_{\text{train}}|}
    \sum_{(a, b) \in \mathcal D_{\text{train}}}
    \log p_W(\,ab \,\big|\, a, b\,).
  \label{eq:bgeo-loss}
\end{equation}
Initialization: each entry of $W$ is drawn i.i.d.\ from
$\mathcal N(0, \sigma^2 / d_{\text{model}})$.

\begin{definition}[Good, bad, and {\color{red}Pareto-optimal} initialization]
\label{def:bgeo-good-init}
Fix thresholds $\tau, \tau_{\text{P}} > 0$.  An initialization
$W^{(0)}$ is
\emph{good} if the gradient flow converges to a $W^{(\infty)}$ whose
recovered $\mathcal K$ satisfies $c^*(W^{(\infty)}) \cdot \Delta_{\min}(\mathcal K) \geq \tau$,
where $c^*$ is the typical per-character amplification (defined
in~\cref{subsec:bgeo-cgap} below).  {\color{red}It is \emph{Pareto-optimal} if the
$\mathcal K$ it reaches lies on the cost-quality Pareto frontier of
bases. \textbf{[The cost-quality frontier rests on the falsified $\Sigma$-orders cost --- see \texttt{cost\_atom.png}.]}}  It is \emph{bad} otherwise.
\end{definition}

Empirically, $\tau \approx 10$ separates the seeds that grok to noise floor
$\sim 10^{-7}$ from the failure modes we have catalogued.

\subsection{Decomposition in Fourier coordinates}
\label{subsec:bgeo-fourier}

The multiplicative-Fourier basis on $\Z/p$ (\cref{subsec:dlog}) gives an
orthonormal basis $\{\chi_k\}_{k=0}^{p-1}$ of $\R^p$.  We decompose the
embedding column-wise:
\begin{equation}
  W \;=\; \sum_{k=0}^{p-1} v_k \,\chi_k^\top,
  \qquad v_k := W\,\chi_k \;\in\; \R^{d_{\text{model}}}.
  \label{eq:bgeo-fourier-decomp}
\end{equation}
Two quantities organize the analysis:
\begin{enumerate}[leftmargin=*]
  \item Per-character energy $e_k(t) := \|v_k(t)\|^2$.
  \item Per-character direction $\hat v_k(t) := v_k(t) / \|v_k(t)\|$.
\end{enumerate}
At random initialization,
\begin{equation}
  e_k^{(0)} \;\overset{\mathrm{i.i.d.}}{\sim}\;
  \frac{\sigma^2}{d_{\text{model}}}\cdot \chi^2_{d_{\text{model}}},
  \qquad \mathbb{E}[e_k^{(0)}] = \sigma^2.
  \label{eq:bgeo-init-stats}
\end{equation}
All characters start with the same expected energy; the seed-to-seed
differences are $O(1/\sqrt{d_{\text{model}}})$ fluctuations.  The
basin-selection question is how these small fluctuations interact with
training dynamics to determine $\mathcal K$.

\subsection{Gradient-flow ODE in Fourier coordinates}
\label{subsec:bgeo-flow}

Projecting \eqref{eq:bgeo-gradflow} onto the Fourier basis:
\begin{equation}
  \frac{\mathrm d v_k}{\mathrm d t} \;=\; -\nabla_{v_k} \mathcal L \;-\; \lambda\, v_k.
  \label{eq:bgeo-vdot}
\end{equation}
Taking the inner product with $2 v_k$ gives the energy dynamics:
\begin{equation}
  \frac{\mathrm d e_k}{\mathrm d t} \;=\; -2\,v_k^\top \nabla_{v_k}\mathcal L \;-\; 2\lambda\, e_k.
  \label{eq:bgeo-edot}
\end{equation}
The naive temptation is to linearize $\nabla_{v_k} \mathcal L$ at small
$v_k$ (i.e., at random init) and extract per-character exponential growth
rates.  We tried this and the prediction failed empirically: the per-character
energies stay roughly \emph{flat} during memorization, not exponentially
growing.  The dynamics that actually matter live at a different timescale and
in a different regime, which we now describe.

\subsection{Equivariance under \texorpdfstring{$(\Z/(p-1))^*$}{(Z/(p-1))*}}
\label{subsec:bgeo-equivariance}

The character group of $\Zp{p-1}$ carries a natural action of $(\Z/(p-1))^*$
by index multiplication: $\chi_k \mapsto \chi_{rk}$ for $r \in (\Z/(p-1))^*$.
This action is realized on tokens by $a \mapsto g^{r \dlog_g a}$, which
preserves the training distribution.  The loss is therefore invariant under
the $(\Z/(p-1))^*$-action on $\WE$, and the gradient flow is equivariant.

\begin{proposition}[Order-class symmetry]
\label{prop:bgeo-order-symmetry}
For any two characters $k, k'$ of the same order $d \mid (p-1)$, the gradient
flow trajectories from $(\Z/(p-1))^*$-related initializations produce
$(\Z/(p-1))^*$-related final states.  In particular, the order multiset of
the final $\mathcal K$ is invariant on each $(\Z/(p-1))^*$-orbit of
initializations.
\end{proposition}

\noindent So $(\Z/(p-1))^*$ collapses 48 character indices into one orbit
within each order class.  But this is only \emph{one} of the loss's
symmetries, and the smaller one.  The empirically dominant symmetry is a
different one.

\subsection{The sign-symmetry of the loss}
\label{subsec:bgeo-sign-symmetry}

For any $(a, b)$, the involution
\begin{equation}
  \sigma : (a, b) \;\longmapsto\; (-a, -b)
  \label{eq:bgeo-sigma}
\end{equation}
preserves the product: $(-a)(-b) = ab \bmod p$.  Hence $\sigma$ preserves the
training data:
\begin{equation}
  (a, b) \in \mathcal D_{\text{train}} \;\iff\; (-a, -b) \in \mathcal D_{\text{train}}
\end{equation}
(under any sign-symmetric sampling, which our split is). The cross-entropy
loss is therefore invariant under the action of $\sigma$ on $\WE$:
\begin{equation}
  \mathcal L(\sigma \cdot W; \theta_{\text{rest}}) \;=\; \mathcal L(W; \sigma \cdot \theta_{\text{rest}}).
\end{equation}
Gradient flow with weight decay respects this symmetry: trajectories
beginning at $\sigma$-related initializations remain $\sigma$-related.

Decompose each character by its parity under~$\sigma$.  A character
$\chi_k$ is \emph{even} if $\chi_k(-a) = \chi_k(a)$ (equivalently,
$\chi_k(-1) = +1$, i.e.\ the order of $\chi_k$ divides $(p-1)/2 = 56$), and
\emph{odd} if $\chi_k(-1) = -1$ (order does not divide $56$, equivalently
the index $k$ has a particular parity in the multiplicative basis).
Equivalently, even characters factor through $\Z/(p-1)/\{\pm 1\} \cong
\Z/((p-1)/2)$; odd characters do not.

Even and odd characters behave very differently under sign-averaging during
training:

\begin{itemize}[leftmargin=*]
  \item \textbf{Even characters} $(\chi(-1) = +1)$: gradient signals from
        $(a,b)$ and from $(-a,-b)$ are \emph{coherent}.  They add
        constructively.  Per-character energy receives sustained, coherent
        amplification.
  \item \textbf{Odd characters} $(\chi(-1) = -1)$: gradient signals from
        $(a,b)$ and from $(-a,-b)$ are \emph{anti-coherent}.  They cancel
        on average.  Per-character energy receives no sustained
        amplification in early training; it changes only via the
        higher-order interaction terms ignored in the linearized picture.
\end{itemize}

This predicts that \emph{even characters get amplified first in early
training, and odd characters lag}.  The orders dividing $56$ for $p = 113$ are
$\{1, 2, 4, 7, 8, 14, 16, 28, 56\}$ --- everything except $112$.  Order-$112$
characters are precisely the odd (primitive) characters.

\subsection{CRT-factor coverage}
\label{subsec:bgeo-crt}

We refine the prediction further.  Recall the CRT decomposition
\begin{equation}
  \Z/112 \;\cong\; \Z/16 \times \Z/7.
\end{equation}
Each character of $\Z/112$ acts on this decomposition.  Indexing characters
by $k \in \{1, \ldots, 111\}$, we classify:
\begin{itemize}[leftmargin=*]
  \item \emph{Pure $\Z/16$-part} ($\chi_k$ trivial on the $\Z/7$ factor):
        $k$ is a multiple of $7$.  Orders in $\{1, 2, 4, 8, 16\}$.
  \item \emph{Pure $\Z/7$-part} ($\chi_k$ trivial on $\Z/16$): $k$ is a
        multiple of $16$.  Orders in $\{1, 7\}$.
  \item \emph{Mixed-factor} (active on both): everything else.  Orders in
        $\{14, 28, 56, 112\}$.
\end{itemize}

The algorithm requires characters that together discriminate elements in
both CRT factors.  Pure-factor characters give weaker overall gradient
signal during memorization than mixed-factor characters, because their
contribution to discriminating $\dlog ab \mod (p-1)$ is restricted to one
CRT factor.

\begin{remark}[Combining the two symmetries]
\label{rem:bgeo-anchor-prediction}
The characters that receive (a) coherent gradient signal under the
sign-symmetry $\sigma$ \emph{and} (b) discriminate both CRT factors are
exactly the \emph{mixed-factor even characters} --- those with order in
$\{14, 28, 56\}$.  We predict these characters dominate the very early
amplification dynamics and are selected first as the basin's ``anchor.''
\end{remark}

\subsection{The anchor commitment phase: empirical findings}
\label{subsec:bgeo-anchor}

We measured the rank of every character in the $W_E$ spectrum at every $50$
epochs of training across $8$ seeds at $d_{\text{model}} = 24$.  For each
seed we identified the \emph{anchor} = the first character to reach rank~$1$
in the spectrum, and recorded its order.

\begin{center}
\begin{tabular}{c|c|c|c}
  seed & anchor character & order & class \\ \hline
  $0$ & $k = 30$ (ep $0$)   & $56$ & mixed \\
  $1$ & $k = 4$ (ep $100$)  & $28$ & mixed \\
  $2$ & $k = 24$ (ep $100$) & $14$ & mixed \\
  $3$ & $k = 2$ (ep $0$)    & $56$ & mixed \\
  $4$ & $k = 46$ (ep $100$) & $56$ & mixed \\
  $5$ & $k = 4$ (ep $0$)    & $28$ & mixed \\
  $6$ & $k = 40$ (ep $50$)  & $14$ & mixed \\
  $7$ & $k = 44$ (ep $100$) & $28$ & mixed \\
\end{tabular}
\end{center}

\textbf{All $8$ seeds anchored on a mixed-factor even character.}  The
distribution: $3$ seeds at order $56$, $3$ at order $28$, $2$ at order $14$.
No seed anchored on a primitive ($112$), a pure-$\Z/16$ character
(orders $2, 4, 8, 16$), or a pure-$\Z/7$ character (order $7$).  Anchor
commitment completed within $\leq 100$ epochs in every case --- in $3$ seeds
the anchor was rank~$1$ already at epoch~$0$ by chance.

This is consistent with \cref{rem:bgeo-anchor-prediction}.  The
intersection of ``even'' (sign-symmetry coherence) and ``mixed-factor''
(both CRT factors active) selects exactly the $26$ characters in orders
$\{14, 28, 56\}$, and the basin's anchor is always one of these.

\subsection{Sequential K recruitment during memorization}
\label{subsec:bgeo-sequential}

After the anchor commits, the rest of $\mathcal K$ is recruited gradually.
Tracking per-MLP-neuron preferred frequencies on seed~$4$, we find:

\begin{itemize}[leftmargin=*]
  \item By epoch $500$ ($\sim 10\%$ of the way to the cliff), $\sim 70\%$
        of MLP neurons already commit to the anchor character $k = 46$.
  \item Through memorization (epochs $500$--$4500$), neurons redistribute
        across the eventual $\mathcal K = \{9, 34, 36, 38, 46\}$.
  \item By the cliff at epoch $\sim 4900$, the final neuron-distribution
        proportions $\{k\!=\!9 \!: 175, k\!=\!34 \!: 107, k\!=\!38 \!: 103,
        k\!=\!46 \!: 127, k\!=\!36 \!: 7\}$ are essentially set.
\end{itemize}

The order in which K characters are recruited respects a pattern: the anchor
(mixed-factor even) commits first; subsequent same-order or
sign-symmetric companions are added next; primitives are added last;
\emph{minor} characters that take only a handful of neurons (like the
order-$28$ character $k = 36$ above with $7$ neurons) often emerge only
near the cliff itself.  In particular, primitive characters that appear in
$\mathcal K$ are recruited \emph{after} the even-character anchor and
companions, never as anchor themselves.

\subsection{The cliff as shortcut decay, not basin selection}
\label{subsec:bgeo-cliff}

During memorization, the model fits the training set via a combination of:
\begin{itemize}[leftmargin=*]
  \item The \emph{algorithm} --- the Fourier-multiplication structure being
        built up across $W_E$, the MLP, and $W_U$ as discussed above.
  \item \emph{Shortcuts} --- high-amplitude features that fit individual
        training examples but do not generalize.
\end{itemize}
Train loss reaches the precision floor during memorization (typically by
epoch $\sim 2000$ in our setup); the loss is jointly minimized by both
components.  The algorithm's contribution to logits is dominated by
shortcuts, so test loss remains at chance.

Weight decay erodes both components, but \emph{at different rates}: the
algorithm has lower norm per training example fit (because one feature
explains many examples efficiently), so it is eroded more slowly than the
shortcuts.  At some critical moment, shortcut amplitudes fall below
algorithm amplitudes in the logit calculation: test loss drops abruptly.
This is the cliff.

The cliff is therefore a \emph{revelation} of the latent algorithm, not its
\emph{discovery}.  The K characters have been preferentially amplified
since the anchor commitment phase; the cliff merely marks the moment when
the algorithm becomes visible at the output.  We have verified this directly:
projecting $W_E$, $W_U$, and the MLP onto K-relevant directions at
mid-memorization reveals a partial algorithm that performs near chance
(strictly better than the original model, which is worse than chance due to
shortcut noise).

\subsection{Within-basin equilibrium: the $c^* \cdot \Delta_{\min}$ relation}
\label{subsec:bgeo-cgap}

Once in the algorithmic basin, the per-character amplification equilibrates
against weight decay.  Define
\begin{equation}
  c_k \;:=\; \|v_k\| \cdot \|u_k\|,
  \qquad
  c^* \;:=\; \operatorname{geom\,mean}\bigl(\{c_k : k \in \mathcal K\}\bigr),
\end{equation}
where $u_k$ is the $\mathcal K$-aligned coefficient in $\WU$.  The product
$c_k$ measures how loudly character $k$ is broadcast into the algorithm's
logit calculation; $c^*$ is the typical per-character amplification across
$\mathcal K$.  The empirical finding across $8$ seeds:

\begin{equation}
  -\log\bigl(\text{test loss}\bigr)
    \;\approx\; 0.46 \cdot c^* \cdot \Delta_{\min}(\mathcal K).
  \label{eq:bgeo-cgap}
\end{equation}

\noindent The fit is sharp: $R^2 \approx 0.99$ over the unsaturated regime
(test loss $\gtrsim 10^{-6}$), and the linear form is uniquely good
compared to alternatives (e.g.\ $c^*$ alone, $\Delta_{\min}$ alone,
$(c^*)^2 \cdot \Delta_{\min}$ all fit substantially worse).

The interpretation is direct: $c^* \cdot \Delta_{\min}$ is the model's
worst-case pre-softmax logit gap.  Cross-entropy on the worst-case
confusable wrong answer is $\sim \exp(-c^* \Delta_{\min})$, so its log is
$-c^* \Delta_{\min}$ up to a constant. The fitted slope of $0.46$ rather
than $1$ reflects calibration of how $\|v_k\|\|u_k\|$ maps to actual logit
scale through the MLP nonlinearity.

\subsection{Weight decay's dual role}
\label{subsec:bgeo-wd-dual}

The empirical $\lambda$-sweep on seed~$4$ shows two distinct regimes:

\begin{itemize}[leftmargin=*]
  \item \textbf{$\lambda \lesssim 0.5$ (too low):} the cliff never fires.
        The memorization basin remains stable; weight decay is too weak to
        erode shortcuts below the algorithm.  Test loss never drops below
        $\sim 10^{-1}$.  This is the failure mode of seed~$5$ at the
        default $\lambda = 1$ as well: its memorization basin is unusually
        robust.
  \item \textbf{$\lambda \in [0.75, 2]$ (working range):} cliff fires, model
        groks.  Within this range, $c^*$ varies modestly with $\lambda$
        and test loss varies with it (factor of $\sim 5\times$ across the
        range, consistent with \eqref{eq:bgeo-cgap}).
\end{itemize}

So $\lambda$ has \emph{two} distinct roles: \emph{(i)}~destabilizing
memorization so the cliff can fire, and \emph{(ii)}~setting the within-basin
amplification equilibrium.  These conflict: $\lambda$ small enough to
allow large $c^*$ may also be too small to destabilize memorization, in
which case the model never grokks.

\subsection{Implications for engineering good initializations}
\label{subsec:bgeo-engineering}

The picture above suggests that ``good initialization'' is determined less by
fine-grained spectral features at epoch $0$ and more by which mixed-factor
even anchor the early dynamics happen to seed the basin with.  Three things
follow:

\begin{enumerate}[leftmargin=*]
  \item The set of accessible $\mathcal K$ is constrained by the anchor: a
        basin seeded with anchor of order $56$ tends to recruit other
        order-$56$ companions; a basin with order-$14$ anchor tends to
        recruit order-$7$ companions.  Different anchors produce
        structurally different final $\mathcal K$.
  \item Engineering a specific final $\mathcal K$ would require steering
        the anchor selection, which in turn requires either biasing the
        early dynamics or using weight-space interventions that survive
        the first few hundred epochs of training.
  \item {\color{red}The Pareto-optimal $\mathcal K$ identified in
        \cref{sec:legendre-channel} may be reachable from a specific
        subset of mixed-factor-even-anchor initializations.  Mapping which
        anchor leads to which Pareto-optimal $\mathcal K$ is the central
        open question. \textbf{[``Pareto-optimal'' uses the falsified $\Sigma$-orders cost; the Legendre being cheap survives, the frontier framing does not --- see \texttt{cost\_atom.png}.]}}
\end{enumerate}

\subsection{Open questions}
\label{subsec:bgeo-open}

\begin{enumerate}[leftmargin=*]
  \item Within mixed-factor even characters, what determines which specific
        one is selected as anchor?  The empirical anchors do not
        correspond simply to the highest-initial-energy character of
        that class.  Some property of the data + MLP init coupling is
        relevant.
  \item Does the picture extend to other primes?  The CRT structure of
        $\Z/(p-1)$ varies with $p$.  The qualitative story (sign-symmetry
        + CRT coverage) should generalize, but the specific anchor-class
        orders depend on $\gcd$ structure.
  \item Can the cliff time be predicted from the anchor-commitment time?
        Empirically, the cliff time across our $8$ seeds varies from
        $\sim 2{,}000$ to $\sim 25{,}000$ epochs, while anchor commitment
        completes by epoch $\sim 100$ in every case.  What sets the
        cliff time?
  \item The bad-initialization region: is it geometrically characterizable
        in $W^{(0)}$-space?  We have classified the failure modes (stuck
        memorization, weak basin, diffuse spectrum) but do not yet know
        which inits flow to which failure.
\end{enumerate}
